\documentclass[11pt]{report}
\usepackage[a4paper,margin=0.9in]{geometry}
\usepackage{amsmath,amssymb}
\usepackage{booktabs}
\usepackage{longtable}
\usepackage{array}
\usepackage[colorlinks=true,linkcolor=blue,urlcolor=blue]{hyperref}
\usepackage{fancyhdr}
\usepackage{xcolor}
\usepackage{titlesec}
\usepackage{enumitem}

\definecolor{hbblue}{HTML}{2A4D8F}
\definecolor{warnred}{HTML}{9C2B2B}
\definecolor{boxgrey}{HTML}{F4F4F6}

\titleformat{\chapter}[display]{\bfseries\Large\color{hbblue}}{\chaptertitlename\ \thechapter}{10pt}{\huge}
\titleformat{\section}{\normalfont\large\bfseries\color{hbblue}}{\thesection}{1em}{}
\pagestyle{fancy}
\fancyhf{}
\fancyhead[R]{\small\itshape Slide Deck Specification}
\fancyfoot[C]{\thepage}

\newsavebox{\kbsave}
\newsavebox{\wbsave}
\newenvironment{keybox}[1]%
{\par\vspace{6pt}\begin{lrbox}{\kbsave}\begin{minipage}{0.88\textwidth}%
{\bfseries\color{hbblue}#1}\par\vspace{4pt}\ignorespaces}%
{\end{minipage}\end{lrbox}%
\begin{center}\setlength{\fboxrule}{1.2pt}\setlength{\fboxsep}{7pt}%
\fcolorbox{hbblue}{boxgrey}{\usebox{\kbsave}}\end{center}\vspace{6pt}\par}
\newenvironment{warnbox}[1]%
{\par\vspace{6pt}\begin{lrbox}{\wbsave}\begin{minipage}{0.88\textwidth}%
{\bfseries\color{warnred}#1}\par\vspace{4pt}\ignorespaces}%
{\end{minipage}\end{lrbox}%
\begin{center}\setlength{\fboxrule}{1.2pt}\setlength{\fboxsep}{7pt}%
\fcolorbox{warnred}{white}{\usebox{\wbsave}}\end{center}\vspace{6pt}\par}

\newcommand{\slide}[2]{\section*{Slide #1 --- #2}\addcontentsline{toc}{section}{Slide #1: #2}}

\title{\textbf{\Huge IYPT Presentation}\\[0.3cm]
\textcolor{hbblue}{\huge Slide-by-Slide Specification}\\[0.5cm]
\large\itshape 12 main slides, 25 reserve. Every figure path verified on disk.}
\author{}
\date{Pinhole Sunglasses --- IYPT 2027 Problem 1}

\begin{document}
\maketitle
\tableofcontents

\chapter{Structure and rules}

\section{The format constraint that drives everything}

IYPT Reporter is 12 minutes. At 45--55 slides that is 13--16 seconds per slide, which is not a
presentation --- it is a slideshow you narrate over while losing the room. The structure here is:

\begin{itemize}[leftmargin=1.2cm]
\item \textbf{12 main slides}, roughly 60\,s each, shown in order.
\item \textbf{25 reserve slides}, shown only when a question demands one.
\end{itemize}

Total is 37 slides --- close to the volume you were planning, but with the depth held in reserve
rather than forced through the clock. Flipping to a prepared backup slide mid-answer reads as
mastery. Running out of time reads as the opposite, regardless of what was on the skipped slides.

\section{Rules for every slide}

\begin{enumerate}[leftmargin=1.2cm]
\item \textbf{One claim per slide.} If you need ``and'' to state the point, it is two slides.
\item \textbf{No paragraph text.} Numbers, one figure, at most three short lines.
\item \textbf{Every number traceable.} If it is on a slide, you can name the CSV it came from.
\item \textbf{Deprecated formulas marked as such} wherever they appear, so nobody thinks you still use them.
\item \textbf{Negatives on your own slides}, in your own words, before anyone asks.
\end{enumerate}

\section{Figures that must be created --- the only real gap}

Everything else is already on disk. Two figures do not exist and cannot be substituted:

\begin{warnbox}{Build these before anything else}
\textbf{FIG-A --- ``One hole per pupil'' overlap diagram.} Hexagonal lattice at 4000\,\textmu m pitch
with a 4000\,\textmu m pupil circle overlaid, showing exactly one hole inside the pupil, plus the
nearest neighbour at 4000\,\textmu m clearly outside. This carries Slide 9, which is the slide that
closes your biggest vulnerability. Without it you are asserting the geometry rather than showing it.

\textbf{FIG-B --- $\mu_g$ saturation histogram.} Distribution of Shen's ghost-membership term across
all 8501 candidates, showing it pinned at one value. This carries Slide 8, your only original
finding against published literature.
\end{warnbox}

Both are quick matplotlib scripts from data you already have. Do them first --- they are the two
slides nobody else at the tournament can show.

\chapter{The 12 main slides}

\slide{1}{The problem, and the tension inside it}

\textbf{Figure:} \texttt{all\_outputs/stage1\_radius\_vs\_diameter.png} --- blur radius vs.\ aperture,
showing both branches and the minimum.

\textbf{On slide:} Title. The figure. Two labels on the curve: ``defocus-limited'' on the falling
branch, ``diffraction-limited'' on the rising branch.

\textbf{Say (45\,s):} ``A short-sighted eye focuses in front of the retina, so every point becomes a
disc. Narrow the aperture and the disc shrinks --- that is why squinting works. But light is a wave,
and a small hole spreads it. Below some size, diffraction grows faster than defocus shrinks. There is
a minimum. Finding it correctly is the whole problem, and `correctly' turns out to be the hard word.''

\slide{2}{The architectural correction}

\textbf{Figure:} Two-panel pipeline diagram --- build in the slide tool, not matplotlib. Top: mask
$\to$ ASM $\to$ pupil $\to$ lens phase $\to$ \textbf{FFT} $\to$ retina, with FFT struck through in
red. Bottom: same chain with \textbf{ASM($v_r$)} in place of the FFT.

\textbf{On slide:} The two chains. One line: $L_{ax} = n'/(P_0-M)$, and $v_r \neq v_f$.

\textbf{Say (90\,s):} ``A lens phase followed by a single FFT gives you the field at the lens's own
back focal plane. For a healthy eye that is the retina and the answer is right. For a myopic eye the
retina is not at the focal plane --- that mismatch is the definition of myopia. So that pipeline
computes a perfectly sharp Airy pattern at the one plane in the system where defocus is exactly zero.
It would model a corrected eye and call it myopic.

We replaced the FFT with a second Angular Spectrum propagation over the real physical distance. That
is not just more accurate --- it makes two formulas unnecessary. Geometric blur is no longer computed
and added; it emerges from propagating a wrongly-curved wavefront the correct distance. And
root-sum-square combination has nothing left to combine.''

\begin{keybox}{This is your strongest slide. Do not rush it.}
Most teams ship this bug and never detect it. Spending 90 seconds here is the single best use of
your twelve minutes.
\end{keybox}

\slide{3}{Validation against the closed form}

\textbf{Figure:} \texttt{data\_collection/stage3\_integrated/plots/validation\_tolerance\_bands.png}

\textbf{On slide:} ``$d_{gap} \to 0$, $M=0$ $\Rightarrow$ must reduce to Fraunhofer''. Three numbers:
$r_{EE}(50\%)$ 2\%, $r_{EE}(90\%)$ 2\%, $r_{EE}(80\%)$ 8\%.

\textbf{Say (45\,s):} ``Setting the gap to zero and myopia to zero, the whole chain must collapse to
the textbook Airy result. It does, to two percent --- except at the 80\% percentile, which needs
eight. That percentile sits almost exactly on the first null, where the encircled-energy curve is
flat, so a tiny error in energy maps to a large error in radius. We checked whether it was
resolution: from grid 1024 to 8192 it never dropped below 2.3\%. It is ill-conditioning, not a bug.''

\begin{keybox}{Volunteer the exception}
Saying the 8\% yourself, with the reason, is worth more than the 2\% figures. It signals you
understand your own error structure rather than reporting only what flatters you.
\end{keybox}

\slide{4}{What the deprecated shortcut costs}

\textbf{Figure:} \texttt{all\_outputs/comparison\_retina\_M4.png} or the RSS-vs-exact comparison plot.

\textbf{On slide:} $r_{tot} \overset{?}{=} \sqrt{r_{geo}^2+r_{diff}^2}$ struck through. Two numbers:
mean divergence 64.8\%, max 192.8\%.

\textbf{Say (45\,s):} ``Root-sum-square assumes two independent Gaussian error terms. Neither blur
here is Gaussian --- defocus gives a roughly uniform disc, diffraction gives an Airy pattern with
rings. Blurs combine by convolving point-spread functions, not by adding radii in quadrature. We did
not assert the shortcut was wrong; we measured it against the exact result. Mean divergence 65\%,
worst case 193\%, and worst at mild myopia --- the regime that matters most.''

\slide{5}{Blur-only optimum}

\textbf{Figure:} \texttt{all\_outputs/stage2\_optimum\_vs\_M\_summary.png}

\textbf{On slide:} The U-curve family. One line: ``optimum 300--400\,\textmu m across 1--8\,D''.

\textbf{Say (40\,s):} ``Minimising blur alone gives 300 to 400 microns, essentially flat across the
whole myopia range. That is a clean result. It is also, as stated, wrong for a real device.''

Deliver that last sentence and pause. It sets up Slide 6.

\slide{6}{The reconciliation --- the pivot of the talk}

\textbf{Figure:} \texttt{brightness\_reconciliation.png}, or
\texttt{stage3a/stage3a\_fig2\_optimum\_vs\_Tmin\_wl550nm\_gap12mm.png}

\textbf{On slide:} $E_{rel}=(\min(a,D_p)/D_p)^2$. Small table: $T_{min}$ 0\% $\to$ 300--400\,\textmu m;
6.25\% $\to$ 1000\,\textmu m; 9\% $\to$ 1250\,\textmu m. Clinical bracket 0.94--1.75\,mm marked.

\textbf{Say (75\,s):} ``Published and commercial pinhole devices sit between 0.94 and 1.75
millimetres. Our blur-only answer is three times smaller. Rather than assume our physics was wrong,
we tested a hypothesis: that brightness is a binding constraint we had omitted.

Illuminance follows from energy conservation --- diffraction redistributes light, it does not destroy
it, so transmission goes as the physical aperture area, not the diffracted spot. Re-running the
optimisation with a floor on retinal illuminance: a floor of six to nine percent, comparable to
ordinary sunglass tinting, moves the optimum into the clinical bracket at every myopia level.

The diffraction physics was never wrong. Blur-only optimisation was answering the wrong question.''

\begin{keybox}{This slide also pre-answers ``which aperture is your answer?''}
Three numbers appear across your work --- 300--400, 900--1250, 530\,\textmu m. They are one U-curve
under three constraint sets. Saying so here means nobody gets to frame it as a contradiction later.
\end{keybox}

\slide{7}{Array physics}

\textbf{Figure:} \texttt{Array\_Design/array\_geometry\_lab/verification\_pdf\_figs/fig\_ghost.png}
alongside \texttt{fig\_viability.png}.

\textbf{On slide:} $d-R_h<R_p$ (viability). $\Delta r = sM/(P_0-M)$ (ghosting), $R^2=0.995$.

\textbf{Say (60\,s):} ``An array adds two questions a single hole does not have. First, does light get
through at all --- pure circle-circle overlap between hole and pupil, and it changes with gaze because
the pupil itself translates as the eye rotates. Second, if two holes are lit at once, each sends a
beam through the lens at a different position, and in a defocused eye they land apart. That
separation is $sM/(P_0-M)$, validated against the full model at $R^2$ of 0.995.

It vanishes at zero myopia by construction --- an in-focus eye converges every beam to one point, so
it cannot ghost. That sanity check is how we caught an earlier version of this formula that was wrong
by two orders of magnitude.''

\slide{8}{Original finding --- a published criterion that does not discriminate}

\textbf{Figure:} \textbf{FIG-B} (must be created). Histogram of $\mu_g$ across 8501 candidates,
pinned at one value.

\textbf{On slide:} ``Shen et al.\ ghost-membership term $\mu_g$'' --- ``8501 candidates evaluated'' ---
``fires on 0''.

\textbf{Say (60\,s):} ``The closest published work to ours scores candidate designs with a
ghost-safety term. We implemented it and ran it across our full candidate set --- 8501 designs
spanning every lattice family, aperture, and pitch we generated. The term saturates. It takes
effectively the same value for every design, so it contributes nothing to the ranking. A criterion
that cannot distinguish between any two designs in a realistic design space is inert.

We are not claiming the paper is wrong. We are reporting that this particular term does not
discriminate at device scale, which matters if anyone intends to use it to select a design.''

\begin{keybox}{Delivery}
Flat, factual, no triumph. The finding is strong enough that overselling it weakens it. This is the
only slide in the deck presenting something nobody has published.
\end{keybox}

\slide{9}{Our design --- and why the uncalibrated parameter does not bind}

\textbf{Figure:} \textbf{FIG-A} (must be created). Lattice at 4000\,\textmu m pitch with 4000\,\textmu m
pupil overlaid, one hole inside, neighbours clearly outside.

\textbf{On slide:} ``pitch 4000\,\textmu m = pupil 4000\,\textmu m''. ``Holes illuminating pupil: 1''.
``Ghost pairs: 0''. ``$\Delta r$ if paired at 8\,D: 615\,\textmu m vs.\ $c_{min}=55$\,\textmu m''.

\textbf{Say (75\,s):} ``The fusion threshold in the published ghost model is a single calibration
constant, and the source paper describes it as qualitative --- it predicts ghost onset at 1.1 dioptres
against a clinical 4. We did not calibrate it. We showed it does not matter for our design.

Our pitch equals the pupil diameter, so only one aperture illuminates the pupil at any gaze angle.
We verified this directly from the generated coordinates: one hole inside the pupil, nearest
neighbour at 4000 microns, well outside. If two were ever lit at 8 dioptres, their images would land
615 microns apart against a 55 micron threshold --- eleven times over, and still over at any
plausible value of that constant.

This design does not survive ghosting by fusing ghost images. It never forms a pair. The parameter is
structurally non-binding here, and that is a geometric argument, not a calibration claim.''

\begin{warnbox}{Immediately follow with the cost --- do not let it be extracted}
``The price is light. One hole per pupil at 400 microns is one percent transmission. That is dark,
and it is the real trade against a dense array, which buys brightness and pays in ghosting it then
has to fuse away. We chose the other side of that trade.''
\end{warnbox}

\slide{10}{Fabrication}

\textbf{Figure:} Photograph of the etched brass array against a scale. If not yet etched, use the
mask SVG rendered, and say so.

\textbf{On slide:} ``350\,\textmu m mask $\to$ 400\,\textmu m etched, 50 holes, 4000\,\textmu m pitch,
0.05\,mm brass''. ``Design $\to$ part traceable row-for-row''.

\textbf{Say (45\,s):} ``Photochemical etching in 0.05 millimetre brass. The etch is isotropic, so it
undercuts the resist and the hole comes out wider than the mask opening --- we compensate by one shim
thickness. The fabrication table traces row-for-row to the design table: the part is the design, not
an approximation to it.''

\slide{11}{Measurement}

\textbf{Figure:} Whatever exists at presentation time. If the logMAR run is done: measured vs.\
predicted. If not: \texttt{ExP3} frosted-screen blur-growth series.

\textbf{On slide (if logMAR done):} Predicted logMAR $+0.33$ at 400\,\textmu m, flat to 0.016 across
0--8\,D. Measured values. Bare-eye baseline, measured.

\textbf{Say (75\,s), flatness version:} ``The model makes one sharp prediction we can test without
knowing anyone's prescription. Through a fixed 400 micron aperture, predicted acuity varies by
0.016 logMAR across eight dioptres of myopia --- a twentieth of one chart line. At that aperture the
aperture sets the acuity, not the refractive error. So subjects with very different uncorrected
vision should read the same line through the same pinhole. That is what we tested.''

\textbf{If no measurement yet:} ``Our single-hole photography shows monotonic blur growth with
defocus, consistent with the model. It is uncalibrated in dioptres, so we present it as
qualitative support, not validation. Calibrated acuity measurement is the next step and we are not
claiming it until we have it.''

\begin{warnbox}{Never state a predicted improvement}
The model produces no bare-eye baseline --- every configuration in the sweep has an aperture. Any
$\Delta$ needs a measured number on both sides.
\end{warnbox}

\slide{12}{Limits}

\textbf{On slide:} Four bullets, no figure.
\begin{itemize}[leftmargin=1cm]
\item Transmission 1\% at 8\,D --- the cost of one-hole-per-pupil.
\item Blind fraction 0.33\% at 8\,D, 0.12\% at 6\,D.
\item $c_{min}$ inherited, argued non-binding, not measured.
\item No calibrated acuity loop closed yet.
\end{itemize}

\textbf{Say (60\,s):} ``What we have not done. Transmission is one percent at high myopia --- dark, and
a genuine cost. At six and eight dioptres there are gaze angles where no hole is lit, 0.12 and 0.33
percent of the time. The fusion constant is inherited and we argued it away geometrically rather than
measuring it. And we have not closed the loop from simulated blur to measured acuity in dioptres.
Those are the four things we would do next.''

\chapter{Reserve deck}

Never shown unless a question calls for one. Group them so you can find one under pressure ---
tab or section markers, not a flat run of 25.

\begin{longtable}{p{0.8cm}p{5.3cm}p{8.2cm}}
\caption{Reserve slides, with the figure each uses.}\\
\toprule
\textbf{\#} & \textbf{Slide} & \textbf{Figure (path under project root)} \\
\midrule
\endhead
\multicolumn{3}{l}{\textbf{Group A --- Metric reliability}}\\
R1 & FWHM ring-hopping & \texttt{data\_collection/stage3\_integrated/plots/retina\_fwhm\_ring\_hopping.png} \\
R2 & Strehl multi-peak investigation & \texttt{data\_collection/stage3\_integrated/plots/strehl\_investigation.png} \\
R3 & On-axis Fresnel closed form, $r=0.997$ & \texttt{all\_outputs/strehl\_vs\_aperture.png} \\
R4 & All EE percentiles & \texttt{data\_collection/stage3\_integrated/plots/ee\_all\_percentiles.png} \\
R5 & MTF curves & \texttt{data\_collection/stage3\_integrated/plots/mtf\_curves.png} \\
R6 & Acuity models compared & \texttt{data\_collection/stage3\_integrated/plots/acuity\_models\_combined.png} \\
R7 & CSF model comparison & \texttt{data\_collection/stage3\_integrated/plots/csf\_comparison.png} \\
\midrule
\multicolumn{3}{l}{\textbf{Group B --- Optics detail}}\\
R8 & Stage-1 heatmap & \texttt{all\_outputs/stage1\_heatmap.png} \\
R9 & Retinal images, M8 & \texttt{all\_outputs/comparison\_retina\_M8.png} \\
R10 & Optimum vs.\ myopia by metric & \texttt{all\_outputs/optimum\_vs\_myopia\_by\_metric.png} \\
R11 & Wavelength sensitivity & \texttt{all\_outputs/stage2\_wavelength\_comparison.png} \\
R12 & Gap sensitivity & \texttt{all\_outputs/stage2\_gap\_comparison.png} \\
R13 & Sensitivity heatmaps & \texttt{data\_collection/stage3\_integrated/plots/sensitivity\_heatmaps\_aperture\_M.png} \\
R14 & Optimum vs.\ Rayleigh criterion & \texttt{all\_outputs/comparison\_optimum\_vs\_rayleigh.png} \\
\midrule
\multicolumn{3}{l}{\textbf{Group C --- Brightness}}\\
R15 & Blur vs.\ brightness trade & \texttt{stage3a/stage3a\_fig1\_blur\_vs\_brightness\_wl550nm\_gap12mm.png} \\
R16 & Optimum vs.\ $T_{min}$ & \texttt{stage3a/stage3a\_fig2\_optimum\_vs\_Tmin\_wl550nm\_gap12mm.png} \\
R17 & Blur-then-constraint & \texttt{stage3a/stage3a\_fig4\_blur\_then\_constraint\_Tmin0.0625.png} \\
R18 & Illuminance in trolands & \texttt{data\_collection/stage3\_integrated/plots/illuminance\_troland.png} \\
R19 & Pareto fronts, all M & \texttt{data\_collection/stage3\_integrated/plots/pareto\_all\_M\_overlay.png} \\
\midrule
\multicolumn{3}{l}{\textbf{Group D --- Array}}\\
R20 & Viability geometry & \texttt{Array\_Design/array\_geometry\_lab/verification\_pdf\_figs/fig\_viability.png} \\
R21 & Packing density & \texttt{Array\_Design/array\_geometry\_lab/verification\_pdf\_figs/fig\_density.png} \\
R22 & Ghosting regime map & \texttt{Array\_Design/ghosting/ghosting\_regime\_map.png} \\
R23 & Joint design space, M8 & \texttt{Array\_Design/collected data/joint\_design\_space\_M8.png} \\
R24 & Gaze-weighted winner, M8 & \texttt{Array\_Design/array\_geometry\_lab/gaze\_weighted\_viz\_figs/fig\_M8.png} \\
R25 & Alignment fractions vs.\ gaze & \texttt{Array\_Design/collected data/alignment\_fractions.png} \\
\bottomrule
\end{longtable}

\begin{keybox}{Two more worth preparing as text-only reserve slides}
\textbf{R26 --- Full per-M winner table} (all eight designs, apertures, margins, blind fractions).\\
\textbf{R27 --- Constants and assumptions} ($P_0=60$\,D, $n'=1.336$, $\lambda=550$\,nm,
$d_{gap}=12$\,mm, $D_{pupil}=4$\,mm, scene luminance 10\,000\,cd/m$^2$). The luminance line matters:
it is the assumption behind every acuity number, and it is the one an examiner is most likely to
probe if you present measured data taken indoors.
\end{keybox}

\chapter{Question preparation}

\begin{longtable}{p{4.6cm}p{9.7cm}}
\caption{The five questions you will get, with the reserve slide to flip to.}\\
\toprule
\textbf{Question} & \textbf{Answer, and where to go} \\
\midrule
\endhead
Why $c_{min}=55$\,\textmu m? &
Shen's, and their own paper calls it a single-parameter qualitative calibration predicting onset at
1.1\,D against clinical 4\,D. We showed it does not bind: pitch equals pupil diameter, no pair forms.
$\to$ Slide 9, then R22. \\
\addlinespace
Which aperture is your answer? &
All three, to three different questions. One U-curve, three constraint sets. $\to$ Slide 6, then R16. \\
\addlinespace
Is FWHM reliable? &
Not per-aperture --- it jumps 4.7 to 154.7\,\textmu m between adjacent apertures. Its coarse trend
tracks clinical data better than $r_{EE}(50\%)$, which is stable but pessimistic by $\sim7\times$. We
optimise on $r_{EE}(50\%)$ and state the bias. $\to$ R1, then R6. \\
\addlinespace
Did you measure that it corrects vision? &
Whatever is true. If not yet: ``No. Our single-hole photography is qualitatively consistent and
uncalibrated. We are not claiming it until we have measured it.'' $\to$ Slide 11. \\
\addlinespace
Why is it so dark? &
One hole per pupil is what makes it ghost-free; 1\% at 8\,D is the price. A dense array buys light
and pays in ghosting it must fuse away. $\to$ Slide 9, then R18. \\
\bottomrule
\end{longtable}

\chapter{Build order}

\begin{enumerate}[leftmargin=1.2cm]
\item \textbf{FIG-A and FIG-B.} Nothing else in the deck is blocked, and these two carry your two
strongest slides. Half a day.
\item \textbf{Slides 2, 6, 8, 9.} The four slides that do the actual work. Build these fully ---
figure, text, and rehearsed script --- before touching anything else.
\item \textbf{Slides 1, 3, 4, 5, 7, 12.} Straightforward: figure already exists, one claim each.
\item \textbf{Reserve deck.} Bulk import; these need no design work, just a figure and a title.
\item \textbf{Slide 10 and 11.} Last, because they depend on fabrication and measurement finishing.
Leave placeholders with the honest fallback text already written.
\item \textbf{Rehearse against the five questions, with a stopwatch.} If any answer runs past 20
seconds or you have to hunt for the reserve slide, the grouping is wrong.
\end{enumerate}

\begin{warnbox}{Time the whole thing before polishing anything}
Twelve minutes is less than it sounds. Run the full deck out loud, with a timer, before you spend a
single hour on visual polish. If it runs long, cut Slide 5 into Slide 6 and shorten Slide 3 --- in
that order. Never cut Slides 2, 6, 8, or 9.
\end{warnbox}

\end{document}
